% Wersja źródłowa: audyt II, jednostka 23
\setcounter{section}{22}
\section{Przestrzenie Montela i nuklearne. Twierdzenie o jądrze Schwartza}
\label{sec:montel-nuklearne23}

\providecommand{\atlasref}[2]{#2}


\noindent\textbf{Wymagane wyniki.} Korzystamy z topologii
przestrzeni funkcji testowych i przestrzeni Schwartza,
twierdzeń Banacha--Steinhausa i Arzelà--Ascoliego,
teorii przestrzeni Hilberta oraz analizy Fouriera.
\medskip

Niech $d\ge1$ będzie liczbą całkowitą, a $\Omega\subseteq\R^d$
niepustym zbiorem otwartym. W poprzednich rozdziałach
zbudowaliśmy przestrzenie
$\mathcal E(\Omega)=C^\infty(\Omega)$,
$\mathcal D(\Omega)$ i $\mathcal S(\R^d)$ wraz z ich
przestrzeniami dualnymi. Obecnie wyjaśnimy, dlaczego w tych
przestrzeniach ograniczoność prowadzi do zwartości oraz w jakim
sensie operator liniowy może być reprezentowany przez jedno jądro
dystrybucyjne. Wszystkie przestrzenie liniowo-topologiczne w tym
rozdziale są nad $\K\in\{\R,\C\}$ i są Hausdorffa. W argumentach
wykorzystujących rozwinięcia Fouriera pracujemy nad $\C$;
wersje rzeczywiste wynikają przez kompleksyfikację.

\subsection{Własność Montela i jej znaczenie}

\noindent\emph{Powiązania tematyczne w Bibliotece matematycznej:} \atlasref{GT-07::thm:topog-zwartosc-sieci}{GT-07: zwartość i podsieci zbieżne}.


\begin{definicja}[Przestrzeń Montela]
\label{def:montel23}
Przestrzenią Montela nazywamy beczkową przestrzeń lokalnie wypukłą
$E$, w której każdy zbiór ograniczony ma zwarte domknięcie.
Jeżeli dodatkowo $E$ jest przestrzenią Frécheta, mówimy o
\emph{przestrzeni Frécheta--Montela}. W tej klasie beczkowość
wynika już z twierdzenia Baire'a, więc wystarczy wymagać
względnej zwartości zbiorów ograniczonych.
\end{definicja}

\begin{twierdzenie}[Klasyczne przestrzenie Montela]
\label{thm:montel-test23}
Dla każdej liczby całkowitej $d\ge1$ i każdego
niepustego zbioru otwartego $\Omega\subseteq\R^d$
przestrzenie $\mathcal E(\Omega)$,
$\mathcal S(\R^d)$ oraz $\mathcal D(\Omega)$, wyposażone
w swoje standardowe topologie lokalnie wypukłe, są przestrzeniami
Montela.
\end{twierdzenie}
\begin{proof}
Dla $\mathcal S$ względną zwartość zbiorów ograniczonych
wykazano w Lemacie~\ref{lem:montel-schwartz22} (dowód
jest taki sam nad $\R$ i nad $\C$). Z Twierdzenia~
\ref{thm:schwartz-frechet17} przestrzeń $\mathcal S$
jest przestrzenią Frécheta, a więc jest beczkowa.

Niech $(K_j)_{j\ge1}$ będzie wyczerpaniem $\Omega$
z Lematu~\ref{lem:exhaustion17}.
Dla ograniczonego $B\subseteq\mathcal E(\Omega)$,
$j\ge1$ i $m\in\N_0$ mamy
$M_{j,m}=\sup_{f\in B}p_{K_j,m}(f)<\infty$,
gdzie $p_{K_j,m}$ jest supremum modułów pochodnych
rzędu co najwyżej $m$ na $K_j$. Możemy założyć,
że $B$ i rozpatrywane $K_j$ są niepuste.
Ponieważ $K_j\subseteq\operatorname{int}K_{j+1}$,
istnieje $\delta_j>0$ takie, że domknięte
$\delta_j$-otoczenie $K_j$ jest zawarte w $K_{j+1}$.
Jeżeli $x,y\in K_j$ i $|x-y|<\delta_j$, odcinek
od $x$ do $y$ leży w $K_{j+1}$. Dla $|\alpha|\le m$
i $f\in B$ z twierdzenia Newtona--Leibniza wynika
\[
 |\partial^\alpha f(x)-\partial^\alpha f(y)|
 \le M_{j+1,m+1}\sum_{i=1}^d|x_i-y_i|.
\]
Rodzina pochodnych jest więc jednostajnie równociągła
i jednostajnie ograniczona na $K_j$, także gdy
$\Omega$ lub $K_j$ nie są wypukłe.
Dla dowolnego ciągu $(f_n)\subseteq B$ twierdzenie
Arzelà--Ascoliego i wybór podciągu przekątniowego
dla przeliczalnej rodziny par $(j,\alpha)$ dają
podciąg $(f_{n_k})$, którego każda pochodna zbiega
jednostajnie na każdym $K_j$. Zatem $(f_{n_k})$
jest ciągiem Cauchy'ego dla wszystkich $p_{K_j,m}$.
Zupełność i metryzowalność z Twierdzenia~
\ref{thm:frechet-cinf17} dają zbieżność w $\mathcal E$.
Dla ciągu $(g_n)$ w domknięciu $B$ wybieramy
$f_n\in B$ w odległości mniejszej niż $1/n$ od $g_n$
w metryce wyznaczającej topologię $\mathcal E$.
Zbieżny podciąg $(f_{n_k})$ daje zbieżny podciąg
$(g_{n_k})$ o tej samej granicy w domknięciu $B$.
Domknięcie jest więc zwarte z kryterium ciągowego
zwartości w przestrzeniach metrycznych.
Przestrzeń $\mathcal E$ jest beczkowa jako przestrzeń Frécheta.

Dla $\mathcal D(\Omega)$ względną zwartość zbiorów
ograniczonych wykazano już w Twierdzeniu~
\ref{thm:D-Montel18}. Pozostaje beczkowość.
Niech $U\subseteq\mathcal D(\Omega)$ będzie beczką,
czyli zbiorem domkniętym, absolutnie wypukłym i pochłaniającym.
Dla każdego etapu $\mathcal D_{K_j}$ przeciwobraz $U$
przez kanoniczne włożenie jest beczką w przestrzeni
Frécheta, a zatem otoczeniem zera w tym etapie.
Funkcjonał Minkowskiego $p_U$ jest seminormą, której
ograniczenia do wszystkich etapów są ciągłe.
Na mocy konstrukcji topologii granicy induktywnej
(Twierdzenie~\ref{thm:LF-universal18}) seminorma
$p_U$ jest ciągła na $\mathcal D(\Omega)$.
Ponieważ $\{f:p_U(f)<1\}\subseteq U$, zbiór $U$
jest otoczeniem zera. Dowodzi to beczkowości,
a tym samym własności Montela.
\end{proof}

\begin{wniosek}[Zbieżność na ograniczonych zbiorach testowych]
\label{cor:montel-dual-seq23}
Niech $E$ będzie przestrzenią Montela, a $(u_n)$ ciągiem w
$E'$ zbieżnym punktowo na $E$ do funkcjonału $u\in E'$.
Wtedy $u_n\to u$ w topologii silnej $\beta(E',E)$.
W szczególności dotyczy to każdej z przestrzeni
z Twierdzenia~\ref{thm:montel-test23}.
\end{wniosek}
\begin{proof}
Ponieważ $E$ jest beczkowa, punktowa ograniczoność
rodziny $\{u_n-u:n\ge1\}$, wynikająca ze zbieżności
punktowej, daje jej równociągłość na mocy twierdzenia
Banacha--Steinhausa. Niech $B\subseteq E$ będzie
niepustym zbiorem ograniczonym i $\varepsilon>0$.
Z równociągłości wybieramy otwarte otoczenie zera $V$
takie, że $|(u_n-u)(v)|<\varepsilon/3$ dla wszystkich
$v\in V$ i $n\ge1$. Zwartość $\overline B$ daje
skończone pokrycie zbiorami $x_i+V$,
$i=1,\ldots,r$, gdzie $x_i\in\overline B$.
Dla dostatecznie dużego $n$ wszystkie
$|(u_n-u)(x_i)|<\varepsilon/3$. Dla $x\in B$
wybieramy $i$ z $x-x_i\in V$; wówczas
$|(u_n-u)(x)|<2\varepsilon/3$.
Stąd $\sup_{x\in B}|(u_n-u)(x)|\le2\varepsilon/3
<\varepsilon$. Dla pustego $B$ warunek jest trywialny.
Jest to zbieżność silna. Odwrotna implikacja wynika
z ograniczoności singletonów.
\end{proof}

\begin{przyklad}[Dlaczego nie wystarczy przestrzeń Banacha]
\label{ex:banach-not-montel23}
Jeżeli $X$ jest nieskończenie wymiarową przestrzenią Banacha,
to jej domknięta kula jednostkowa nie jest zwarta w topologii
normowej, na mocy lematu Riesza. Dlatego taka przestrzeń
nie jest przestrzenią Montela. Własność Montela
jest tu skutkiem szczególnie mocnej topologii przestrzeni
funkcji gładkich, a nie samej zupełności.
\end{przyklad}

\subsection{Operatory nuklearne w przestrzeniach Banacha}

\begin{definicja}[Operator nuklearny]
\label{def:nuclear-banach23}
Niech $X,Y$ będą przestrzeniami Banacha. Operator liniowy
$T:X\to Y$ nazywamy \emph{nuklearnym}, jeżeli istnieją
$\ell_j\in X'$ i $y_j\in Y$, $j\ge1$, takie, że
\[
 Tx=\sum_{j=1}^\infty \ell_j(x)y_j\quad(x\in X),
 \qquad \sum_{j=1}^\infty\|\ell_j\|\,\|y_j\|<\infty.
\]
Normą nuklearną $\|T\|_{\mathcal N}$ nazywamy infimum
ostatniej sumy po wszystkich takich przedstawieniach.
Przedstawienie dopuszcza także skończenie wiele składników.
Dla ustalonego $x\in X$ szereg jest bezwzględnie
zbieżny w normie $Y$, ponieważ
$\sum_j\|\ell_j(x)y_j\|\le\|x\|\sum_j\|\ell_j\|\|y_j\|$,
a $Y$ jest zupełna.
\end{definicja}

\begin{twierdzenie}[Nuklearność implikuje zwartość]
\label{thm:nuclear-compact23}
Jeśli $X,Y$ są przestrzeniami Banacha i $T:X\to Y$
jest operatorem nuklearnym, to $T$ jest ograniczony i zwarty.
\end{twierdzenie}
\begin{proof}
Dla przedstawienia z Definicji~\ref{def:nuclear-banach23}
niech $T_Nx=\sum_{j=1}^N\ell_j(x)y_j$.
Dla każdego $x\in X$ mamy
$\|Tx\|\le\|x\|\sum_j\|\ell_j\|\|y_j\|$,
więc $T$ jest ograniczony. Każdy $T_N$ ma skończony
rząd i jest ograniczony. Ponadto
$\|T-T_N\|\le\sum_{j>N}\|\ell_j\|\|y_j\|\to0$.
Przeto $T$ jest ograniczony i jest granicą operatorów
skończonego rzędu w normie operatorowej. Granica
operatorów zwartych w normie operatorowej jest zwarta:
obraz kuli jednostkowej można dla dowolnego $\varepsilon>0$
pokryć skończoną liczbą kul o promieniu $\varepsilon$
za pomocą dostatecznie dobrego przybliżenia $T_N$.
Dokładniej, wybieramy $N$ z $\|T-T_N\|<\varepsilon/2$.
Obraz kuli jednostkowej przez $T_N$ jest ograniczony
w przestrzeni skończenie wymiarowej, więc ma skończone
pokrycie kulami o promieniu $\varepsilon/2$.
Te same środki z promieniem $\varepsilon$ pokrywają
obraz przez $T$. Zupełność $Y$ i całkowita ograniczoność
obrazu zapewniają zwartość jego domknięcia.
\end{proof}

\begin{stwierdzenie}[Norma nuklearna i składanie operatorów]
\label{prop:nuclear-ideal23}
Niech $X,Y$ będą przestrzeniami Banacha nad $\K$.
Operatory nuklearne $X\to Y$ tworzą przestrzeń liniową,
a $\|\cdot\|_{\mathcal N}$ jest na niej normą,
spełniającą $\|T\|\le\|T\|_{\mathcal N}$.
Jeżeli dodatkowo $X_0,Y_0$ są przestrzeniami Banacha
i $A:X_0\to X$, $B:Y\to Y_0$ są ograniczonymi
operatorami liniowymi, to $BTA$ jest nuklearny oraz
$\|BTA\|_{\mathcal N}\le\|B\|\|T\|_{\mathcal N}\|A\|$.
\end{stwierdzenie}
\begin{proof}
Przedstawienia nuklearne dwóch operatorów można połączyć
w jeden szereg dla ich sumy; suma norm składników jest
sumą odpowiednich oszacowań. Skalowanie jednego z czynników
każdego składnika przez skalar daje przedstawienie
dla wielokrotności operatora. Biorąc infima, otrzymujemy
nierówność trójkąta i jednorodność (dla niezerowego
skalaru stosujemy też skalowanie przez jego odwrotność).
Oszacowanie $\|T\|\le\sum_j\|\ell_j\|\|y_j\|$
z poprzedniego dowodu daje $\|T\|\le\|T\|_{\mathcal N}$,
więc norma nuklearna oddziela punkty.
Dla złożenia mamy
$BTAx=\sum_j(\ell_j\circ A)(x)By_j$ oraz
\[
 \sum_j\|\ell_j\circ A\|\|By_j\|
 \le\|B\|\|A\|\sum_j\|\ell_j\|\|y_j\|.
\]
Ciągłość $B$ uzasadnia zastosowanie go do sumy szeregu,
a wzięcie infimum daje żądane oszacowanie.
\end{proof}

\begin{przyklad}[Operator zwarty, który nie jest nuklearny]
\label{ex:compact-nonnuclear23}
Na przestrzeni Hilberta $\ell^2(\N;\K)$ z kanoniczną
bazą ortonormalną $(e_n)_{n\ge1}$ rozważmy operator
$D(x_n)_{n\ge1}=(x_n/n)_{n\ge1}$. Jest liniowy i
ograniczony, gdyż $\|Dx\|_2\le\|x\|_2$.
Iloczyn skalarny przyjmujemy liniowy w pierwszym argumencie.
Operator $D$ jest zwarty, gdyż operatory skończonego rzędu
$D_N=\sum_{n\le N}n^{-1}\langle\cdot,e_n\rangle e_n$
spełniają $\|D-D_N\|=1/(N+1)$:
oszacowanie z góry wynika z sumy kwadratów ogona,
a równość uzyskujemy dla $x=e_{N+1}$. Nie jest jednak
nuklearny. Istotnie, dla każdego przedstawienia nuklearnego
$D=\sum_j\ell_j(\cdot)y_j$ i każdego $N$ mamy
\[
\sum_{n=1}^N|\langle De_n,e_n\rangle|
\le\sum_j
\Big(\sum_{n=1}^N|\ell_j(e_n)|^2\Big)^{1/2}
\Big(\sum_{n=1}^N|\langle y_j,e_n\rangle|^2\Big)^{1/2}
\le\sum_j\|\ell_j\|\|y_j\|.
\]
Istotnie, w pierwszej nierówności stosujemy nierówność
trójkąta do szeregu dla każdego $n$, zamieniamy sumę
skończoną z szeregiem nieujemnym i stosujemy nierówność
Cauchy'ego--Schwarza do sumy po $n$. W drugiej nierówności
korzystamy z twierdzenia Riesza o reprezentacji funkcjonałów
na przestrzeni Hilberta oraz z nierówności Bessela:
$\sum_{n=1}^N|\ell_j(e_n)|^2\le\|\ell_j\|^2$ i
$\sum_{n=1}^N|\langle y_j,e_n\rangle|^2\le\|y_j\|^2$.
Lewa strona jest równa $\sum_{n=1}^N1/n$ i rośnie
nieograniczenie. Zatem zwartość jest własnością ściśle
słabszą niż nuklearność.
\end{przyklad}

\subsection{Nuklearność przestrzeni lokalnie wypukłych}

\begin{definicja}[Przestrzeń nuklearna]
\label{def:nuclear-lcs23}
Niech $E$ będzie przestrzenią lokalnie wypukłą Hausdorffa,
$\mathcal P$ rodziną wszystkich ciągłych seminorm na $E$.
Dla $p\in\mathcal P$ definiujemy przestrzeń Banacha $E_p$
jako uzupełnienie przestrzeni unormowanej $E/\ker p$
w normie $\|x+\ker p\|=p(x)$.
Przez $[x]_p$ oznaczamy obraz $x\in E$ w $E_p$.
Jeżeli $p\le q$, to $\ker q\subseteq\ker p$ i
odwzorowanie $[x]_q\mapsto[x]_p$ ma normę co najwyżej $1$,
więc jednoznacznie przedłuża się do ograniczonego operatora
$j_{q,p}:E_q\to E_p$ o gęstym obrazie.
Dla $p\le q\le r$ zachodzi $j_{r,p}=j_{q,p}j_{r,q}$:
obie strony zgadzają się na gęstym obrazie $E$ w $E_r$.
Przestrzeń $E$ nazywamy \emph{nuklearną}, gdy dla każdego
$p\in\mathcal P$ istnieje $q\in\mathcal P$, $q\ge p$,
takie, że $j_{q,p}$ jest operatorem nuklearnym.
\end{definicja}

\begin{lemat}[Kryterium dla rodziny seminorm wyznaczającej topologię]
\label{lem:nuclear-cofinal23}
Niech $E$ będzie lokalnie wypukłą przestrzenią Hausdorffa,
a $\mathcal P_0$ rodziną ciągłych seminorm wyznaczającą
jej topologię i skierowaną ku górze względem nierówności.
Jeżeli dla każdego $p_0\in\mathcal P_0$ istnieje
$q_0\in\mathcal P_0$ z $q_0\ge p_0$ i nuklearnym
$j_{q_0,p_0}$, to $E$ jest nuklearna.
\end{lemat}
\begin{proof}
Niech $p$ będzie dowolną ciągłą seminormą. Z ciągłości
i skierowania rodziny wynika $p\le C p_0$ dla pewnych
$C\ge1$ i $p_0\in\mathcal P_0$: odpowiednie skończone
przecięcie kul seminormowych zastępujemy jedną kulą
i stosujemy skalowanie. Wybierzmy $q_0$ jak w założeniu
i połóżmy $q=Cq_0$. Wtedy $q\ge p$.
Odwzorowanie $E_q\to E_{q_0}$ dane przez identyczność
na klasach jest ograniczone, podobnie jak
$E_{p_0}\to E_p$, którego norma nie przekracza $C$.
Ich złożenie z $j_{q_0,p_0}$ jest nuklearne na mocy
Stwierdzenia~\ref{prop:nuclear-ideal23} i zgadza się
z $j_{q,p}$ na gęstym obrazie $E$ w $E_q$.
\end{proof}

\begin{uwaga}[Ważne rozróżnienie]
\label{rem:nuclear-different23}
Nuklearność przestrzeni lokalnie wypukłej jest własnością
\emph{całej rodziny seminorm}, a nie twierdzeniem,
że identyczność na tej przestrzeni ma jedno nuklearne
przedstawienie. Jeśli przestrzeń Banacha $X$ jest nuklearna
w sensie Definicji~\ref{def:nuclear-lcs23}, to $X$
ma skończony wymiar: dla $p(x)=\|x\|$ nuklearność
daje ciągłą seminormę $q\ge p$. Z ciągłości $q$
w topologii normowej wynika $q\le Cp$, więc $q$
jest normą równoważną $p$ i oba uzupełnienia
identyfikujemy z $X$. Nuklearne $j_{q,p}$ po
złożeniu z ograniczoną identycznością $X_p\to X_q$
daje nuklearną identyczność na $X$, a ta jest zwarta
na mocy Twierdzenia~\ref{thm:nuclear-compact23}.
Lemat Riesza wyklucza nieskończony wymiar.
Odwrotnie, w skończenie wymiarowej przestrzeni Banacha
każde $j_{q,p}$ ma skończony rząd i jest nuklearne
(przedstawienie w bazie obrazu z ciągłymi współrzędnymi);
można wziąć $q=p$.
\end{uwaga}

\begin{twierdzenie}[Nuklearność przestrzeni szybko malejących ciągów]
\label{thm:nuclear-s23}
Niech $d\ge1$ będzie liczbą całkowitą, $I=\N_0^d$
oraz $a_\nu\in\K$. Definiujemy
\[
 s(I)=\bigl\{a=(a_\nu)_{\nu\in I}:
 p_k(a)=\sup_{\nu\in I}(1+|\nu|)^k|a_\nu|<\infty
 \text{ dla każdego }k\in\N_0\bigr\},
 \qquad |\nu|=\nu_1+\cdots+\nu_d.
\]
Przestrzeń $s(I)$ z topologią seminorm $(p_k)$ jest
nuklearną przestrzenią Frécheta--Montela.
\end{twierdzenie}
\begin{proof}
Normy $p_k$ są rosnące i oddzielają punkty, więc
wyznaczają metryzowalną topologię lokalnie wypukłą.
Niech $(a^{(j)})$ będzie ciągiem Cauchy'ego w każdej
normie $p_k$. Granice współrzędnych
$a_\nu=\lim_j a_\nu^{(j)}$ istnieją dzięki zupełności $\K$.
Dla ustalonego $k$ i $\varepsilon>0$ istnieje $J$
takie, że dla $i,j\ge J$ i każdego $\nu$ zachodzi
$(1+|\nu|)^k|a_\nu^{(i)}-a_\nu^{(j)}|<\varepsilon$.
Przechodząc z $j$ do nieskończoności, otrzymujemy
$p_k(a^{(i)}-a)\le\varepsilon$ dla $i\ge J$.
W szczególności $p_k(a)<\infty$, więc $a\in s(I)$
i $a^{(i)}\to a$ w jego topologii. Jest to przestrzeń
Frécheta.

Dla $m\in\N_0$ oznaczmy $E_m=s(I)_{p_m}$.
Jest to przestrzeń
\[
 c_{0,m}(I)=\{a:(1+|\nu|)^m a_\nu\to0
 \text{ poza skończonymi podzbiorami }I\}
\]
z normą $p_m$. Mnożenie współrzędnych przez
$(1+|\nu|)^m$ jest izometrią tej przestrzeni na
zupełną przestrzeń $c_0(I;\K)$. Ponieważ
$(1+|\nu|)^m|a_\nu|\le p_{m+1}(a)/(1+|\nu|)$,
mamy $s(I)\subseteq c_{0,m}(I)$.
Ciągi o skończonym nośniku należą do $s(I)$
i są gęste w $c_{0,m}(I)$ przez obcinanie ogonów.
To uzasadnia identyfikację uzupełnienia.

Ustalmy $k$ i wybierzmy całkowite $m>k+d$.
Niech $e_\nu$ będzie ciągiem równym $1$ na indeksie $\nu$
i $0$ na pozostałych. Funkcjonał
$\mathrm{ev}_\nu(a)=a_\nu$ na $E_m$ ma normę
$(1+|\nu|)^{-m}$ (równość osiąga odpowiednio
przeskalowany $e_\nu$), a $p_k(e_\nu)=(1+|\nu|)^k$.
Zatem
\[
 \sum_{\nu\in I}\|\mathrm{ev}_\nu\|_{E_m'}
 \|e_\nu\|_{E_k}
 =\sum_{\nu\in I}(1+|\nu|)^{k-m}<\infty.
\]
Istotnie, liczba rozwiązań równania
$\nu_1+\cdots+\nu_d=r$ w $\N_0^d$ wynosi
$\binom{r+d-1}{d-1}$ (rozmieszczamy $d-1$
przegród między $r$ jednostkami).
Jest ograniczona przez $C_d(1+r)^{d-1}$.
Otrzymujemy szereg o wyrazach ograniczonych przez
$C_d(1+r)^{k-m+d-1}$, zbieżny, gdyż
$k-m+d-1<-1$. Szereg
$j_{p_m,p_k}a=\sum_{\nu\in I}a_\nu e_\nu$
jest zbieżny w normie $E_k$ i ma podane współrzędne.
Jest to przedstawienie nuklearne przejścia kanonicznego.
Lemat~\ref{lem:nuclear-cofinal23} daje nuklearność $s(I)$.

Niech $B\subseteq s(I)$ będzie niepustym zbiorem
ograniczonym i $M_k=\sup_{a\in B}p_k(a)$.
Dla każdego $R\ge0$ mamy
\[
 \sup_{a\in B,\,|\nu|>R}(1+|\nu|)^k|a_\nu|
 \le \frac{M_{k+1}}{1+R}.
\]
Dla dowolnego ciągu w $B$ wybieramy podciąg
przekątniowy zbieżny na każdej współrzędnej
(korzystając z ograniczoności w $p_0$).
Na skończonym zbiorze $|\nu|\le R$ jest on
ciągiem Cauchy'ego w każdej normie ważonej.
Na pozostałych indeksach różnica dwóch wyrazów
ma normę ważoną ograniczoną przez $2M_{k+1}/(1+R)$.
Wybierając najpierw $R$, a potem dostatecznie późne
wyrazy, otrzymujemy ciąg Cauchy'ego w każdej normie.
Z zupełności wynika jego zbieżność w $s(I)$.
Jak w dowodzie Twierdzenia~\ref{thm:montel-test23},
metryzowalność pozwala przejść do zwartości domknięcia
$B$. Beczkowość wynika z własności Frécheta.
\end{proof}

\begin{twierdzenie}[Nuklearność $C^\infty(\Omega)$]
\label{thm:nuclear-cinfty23}
Dla każdego niepustego zbioru otwartego $\Omega\subseteq\R^d$
przestrzeń $\mathcal E(\Omega)=C^\infty(\Omega)$
ze standardową topologią Frécheta jest nuklearna.
\end{twierdzenie}
\begin{proof}
Rozpatrzmy najpierw przypadek zespolony. Ustalmy
niepusty zwarty $K\Subset\Omega$ i $m\in\N_0$.
Niech $p_{K,m}(f)=\max_{|\alpha|\le m}\sup_{x\in K}
|\partial^\alpha f(x)|$. Z Lematu~\ref{lem:cutoffs17}
wybieramy $\chi\in\mathcal D(\Omega)$ równą $1$
w otoczeniu $K$, o nośniku $L=\supp\chi\Subset\Omega$.
W szczególności $K\subseteq L$. Funkcję $\chi f$
przedłużamy zerem na $\R^d$, co daje funkcję gładką,
ponieważ $L\Subset\Omega$.
Wybieramy $R>0$ z $L\subset(-R,R)^d$ i przedłużamy
tę funkcję okresowo z sześcianu $Q=[-R,R]^d$.
Przedłużenie jest gładkie, ponieważ funkcja znika
w otoczeniu brzegu sześcianu.
Przy znormalizowanej mierze $(2R)^{-d}\lambda_d$
na $Q$ funkcje $e_n(x)=e^{i\pi n\cdot x/R}$,
$n\in\mathbb Z^d$, tworzą bazę ortonormalną w $L^2$
torusa. Niech
\[
 c_n(f)=(2R)^{-d}\int_Q(\chi f)(x)
 e^{-i\pi n\cdot x/R}\,\mathrm d\lambda_d(x).
\]
Ponieważ $(1-\Delta)e_n=(1+\pi^2|n|^2/R^2)e_n$,
dla każdego $r\in\N_0$ całkowanie przez części
bez wyrazów brzegowych daje
\[
 (1+\pi^2|n|^2/R^2)^r c_n(f)
 =(2R)^{-d}\int_Q (1-\Delta)^r(\chi f)(x)
 \overline{e_n(x)}\,\mathrm d\lambda_d(x).
\]
Z reguły Leibniza i zwartości $L$ otrzymujemy
\[
 |c_n(f)|\le C_r(1+|n|)^{-2r}p_{L,2r}(f).
\]
Stała zależy od $r,\chi,R,d$, lecz nie od $f,n$.
W szczególności $c_n$ znika na $\ker p_{L,2r}$
i przedłuża się do funkcjonału na uzupełnieniu
$E_{p_{L,2r}}$ z normą ograniczoną przez
$C_r(1+|n|)^{-2r}$.

Niech $v_n=[\chi e_n]_{p_{K,m}}$.
Ponieważ $\chi=1$ w otoczeniu $K$, klasa ta ma
na $K$ wszystkie pochodne do rzędu $m$ takie jak $e_n$;
jej norma nie przekracza $C_{R,m}(1+|n|)^m$.
Wybierzmy $r$ tak, by $2r>m+d$.
Szereg $\sum_n (1+|n|)^{m-2r}$ jest zbieżny:
liczba punktów $\mathbb Z^d$ z normą maksimum równą
$s\ge1$ wynosi $(2s+1)^d-(2s-1)^d
\le C_d(1+s)^{d-1}$, a normy euklidesowa i maksimum
są równoważne.

Rozwinięcie $\sum_n c_n(f)e_n$ jest jednostajnie
zbieżne na $Q$ wraz ze wszystkimi pochodnymi do rzędu
$m$, na mocy powyższych oszacowań.
Sumę szeregu bez pochodnych identyfikujemy z okresowym
przedłużeniem $\chi f$: sumy częściowe zbiegają
w $L^2(Q)$ do $\chi f$ z zupełności układu
trygonometrycznego, a jednostajna zbieżność daje
ten sam limit w $L^2(Q)$. Obie funkcje są ciągłe,
więc równość prawie wszędzie daje równość wszędzie.
Jednostajna zbieżność pochodnych pozwala różniczkować
szereg do rzędu $m$ (stosujemy twierdzenie
Newtona--Leibniza kolejno do każdej współrzędnej).
Ponieważ $\chi f=f$ w otoczeniu $K$ i $\chi e_n=e_n$
w tym samym otoczeniu, otrzymujemy w uzupełnieniu
\[
 [f]_{p_{K,m}}=\sum_{n\in\mathbb Z^d}c_n(f)v_n.
\]
Nie chodzi jedynie o wartości funkcji na $K$,
lecz także o pochodne do rzędu $m$.
Ponadto
\[
 \sum_n\|c_n\|_{(E_{p_{L,2r}})'}\|v_n\|_{E_{p_{K,m}}}
 \le C\sum_n(1+|n|)^{m-2r}<\infty.
\]
Ponieważ $p_{K,m}\le p_{L,2r}$, szereg definiuje
operator nuklearny $E_{p_{L,2r}}\to E_{p_{K,m}}$.
Zgadza się on z przejściem kanonicznym na gęstym
obrazie $E$, a zatem na całym uzupełnieniu.

W przypadku rzeczywistym stosujemy tę samą konstrukcję
do rzeczywistego $f$, ale bierzemy część rzeczywistą
szeregu. Składnik zastępujemy sumą
$\operatorname{Re}c_n(f)\operatorname{Re}v_n-
\operatorname{Im}c_n(f)\operatorname{Im}v_n$,
gdzie części rzeczywiste i urojone klas reprezentowane
są przez odpowiednie części funkcji $\chi e_n$
(możemy wybrać $\chi$ rzeczywistą).
Normy obu rzeczywistych funkcjonałów i obu wektorów
nie przekraczają odpowiadających oszacowań zespolonych,
więc szereg nuklearny ma sumę norm ograniczoną przez
dwukrotność poprzedniego oszacowania.

Standardowe seminormy $p_{K,m}$ tworzą rodzinę
skierowaną ku górze: dla dwóch zbiorów bierzemy
ich sumę mnogościową, a dla rzędów maksimum.
Wyznaczają topologię $\mathcal E(\Omega)$.
Lemat~\ref{lem:nuclear-cofinal23} kończy dowód
w obu przypadkach.
\end{proof}

\begin{uwaga}[Hermite i przestrzeń Schwartza]
\label{rem:hermite-nuclear23}
Przestrzeń Schwartza $\mathcal S(\R^d)$ jest nuklearna:
współczynniki w bazie funkcji Hermite'a wyznaczają
izomorfizm liniowo-topologiczny
$\mathcal S(\R^d)\simeq s(\N_0^d)$.
Dokładne oszacowania i dowód tego izomorfizmu podano
w Uzupełnieniu~\ref{sec:hermite-supp23}.
Nuklearność $\mathcal D_K(\Omega)$ oraz
$\mathcal D(\Omega)$ zostanie poniżej wyprowadzona
z udowodnionych twierdzeń trwałości; nie jest to
formalna konsekwencja samej nuklearności $\mathcal S$.
\end{uwaga}

\subsection{Kryterium hilbertowskie i trwałość nuklearności}
\label{subsec:nuclear-permanence23}

W zastosowaniach do funkcji testowych potrzebujemy dwóch operacji:
przejścia do podprzestrzeni i do granicy induktywnej. Dowiedziemy
odpowiednich twierdzeń, wychodząc wyłącznie z
Definicji~\ref{def:nuclear-lcs23}. Ciągłą seminormę $h$ na przestrzeni lokalnie wypukłej $E$
nazywamy \emph{hilbertowską}, jeśli istnieje dodatnia
półokreślona hermitowska forma półtoraliniowa $b$ (nad $\R$:
symetryczna forma dwuliniowa), taka że $h(x)^2=b(x,x)$
dla każdego $x\in E$. Przyjmujemy liniowość w pierwszym
argumencie. Nierówność Cauchy'ego--Schwarza dla $b$
zapewnia $b(x,y)=0$ dla $x\in\ker h$ i dowolnego $y$,
więc $b$ wyznacza iloczyn skalarny na $E/\ker h$.
Przez $H_h=E_h$ oznaczamy jego uzupełnienie hilbertowskie.
Dla $h\le r$ operator $j_{r,h}:H_r\to H_h$ jest
przejściem kanonicznym z Definicji~\ref{def:nuclear-lcs23}.
Nie wymagamy, by $E$ była metryzowalna.

Przypomnijmy, że ograniczony operator liniowy $J:H\to G$
między przestrzeniami Hilberta jest operatorem
Hilberta--Schmidta, gdy $\sum_{i\in I}\|Ju_i\|^2<\infty$
dla bazy ortonormalnej $(u_i)_{i\in I}$ przestrzeni $H$.
Dla nieujemnych składników sumę po dowolnym zbiorze
indeksów definiujemy jako supremum sum skończonych.
Jeśli $(v_j)_{j\in J_0}$ jest bazą ortonormalną $G$,
to z tożsamości Parsevala i zamiany sum nieujemnych
otrzymujemy
\[
 \sum_{i\in I}\|Ju_i\|^2
 =\sum_{i\in I}\sum_{j\in J_0}|\langle Ju_i,v_j\rangle|^2
 =\sum_{j\in J_0}\|J^*v_j\|^2.
\]
Dowodzi to niezależności od wyboru obu baz oraz
$\|J^*\|_{\mathrm{HS}}=\|J\|_{\mathrm{HS}}$, gdzie
$\|J\|_{\mathrm{HS}}^2=\sum_i\|Ju_i\|^2$.

\begin{lemat}[Złożenie operatorów Hilberta--Schmidta]
\label{lem:hs-compose-nuclear23}
Niech $H_2,H_1,H_0$ będą przestrzeniami Hilberta nad $\K$, a
$A:H_2\to H_1$ oraz $B:H_1\to H_0$ operatorami
Hilberta--Schmidta. Wówczas $BA:H_2\to H_0$ jest
operatorem nuklearnym w sensie
Definicji~\ref{def:nuclear-banach23}, przy czym
$\|BA\|_{\mathcal N}\le\|A\|_{\mathrm{HS}}\|B\|_{\mathrm{HS}}$.
\end{lemat}
\begin{proof}
Niech $(e_j)_{j\in J}$ będzie dowolną bazą ortonormalną $H_1$.
Dla każdego $x\in H_2$ zachodzi
$BAx=\sum_{j\in J}\langle Ax,e_j\rangle Be_j$,
gdyż rozwinięcie $Ax$ w bazie jest zbieżne w $H_1$,
a $B$ jest ograniczony. Ponadto
\[
 \sum_{j\in J}\|A^*e_j\|\,\|Be_j\|
 \le\Bigl(\sum_{j\in J}\|A^*e_j\|^2\Bigr)^{1/2}
     \Bigl(\sum_{j\in J}\|Be_j\|^2\Bigr)^{1/2}
 =\|A\|_{\mathrm{HS}}\|B\|_{\mathrm{HS}}.
\]
Przy sumowaniu po dowolnym zbiorze indeksów rozumiemy supremum
sum po jego skończonych podzbiorach. Skończoność powyższych sum
kwadratów zapewnia, że zbiory indeksów z
$\|A^*e_j\|>0$ lub $\|Be_j\|>0$ są co najwyżej przeliczalne:
dla każdego $k\ge1$ tylko skończenie wiele składników
sumy kwadratów może przekraczać $1/k$.
Pozostałe składniki przedstawienia są zerowe,
nawet gdy $H_1$ nie jest ośrodkowa. Funkcjonały
$x\mapsto\langle Ax,e_j\rangle$ mają normy $\|A^*e_j\|$,
a wektory $Be_j$ normy $\|Be_j\|$. Otrzymujemy zatem
przedstawienie nuklearne $BA$ z wymaganym oszacowaniem.
\end{proof}

\begin{twierdzenie}[Hilbertowskie kryterium nuklearności]
\label{thm:hilbert-nuclear23}
Niech $E$ będzie przestrzenią lokalnie wypukłą Hausdorffa
nad $\K$. Następujące warunki są równoważne:
\begin{enumerate}[label=\textup{(\roman*)},leftmargin=*]
\item $E$ jest nuklearna w sensie
Definicji~\ref{def:nuclear-lcs23};
\item każda ciągła seminorma $p$ na $E$ jest ograniczona
z góry przez pewną ciągłą seminormę hilbertowską $h$,
a dla każdej ciągłej seminormy hilbertowskiej $h$ istnieje
ciągła seminorma hilbertowska $r\ge h$, dla której
kanoniczny operator $H_r\to H_h$ jest operatorem
Hilberta--Schmidta.
\end{enumerate}
\end{twierdzenie}
\begin{proof}
\emph{(i)$\Rightarrow$(ii).} Ustalmy dowolną ciągłą seminormę $p$.
Z nuklearności wybieramy $q\ge p$ i przedstawienie operatora
$j_{q,p}:E_q\to E_p$ w postaci
\[
 j_{q,p}z=\sum_{j\ge1}a_j\ell_j(z)y_j,\qquad
 A=\sum_{j\ge1}|a_j|<\infty,\quad
 \|\ell_j\|\le1,\quad\|y_j\|\le1.
\]
Takie przedstawienie uzyskujemy, usuwając składniki
zerowe i dzieląc każdy funkcjonał oraz każdy wektor
przez jego normę, a iloczyn norm przenosząc do
współczynnika $a_j$. Można więc nawet przyjąć $a_j\ge0$. Jeśli $A=0$, to $p=0$ i wystarcza
seminorma zerowa. W przeciwnym razie kładziemy
$h(x)^2=A\sum_j|a_j|\,|\ell_j([x]_q)|^2$.
Jest to seminorma hilbertowska zadana przez formę
\[
 b(x,y)=A\sum_j|a_j|\ell_j([x]_q)
 \overline{\ell_j([y]_q)}.
\]
W przypadku rzeczywistym pomijamy sprzężenie.
Szereg jest bezwzględnie zbieżny, ponieważ
$|\ell_j([x]_q)|\le q(x)$ i $\sum_j|a_j|=A$.
Forma jest dodatnia półokreślona, a $h(x)\le Aq(x)$
dowodzi ciągłości seminormy.
Z nierówności Cauchy'ego--Schwarza wynika
\[
 p(x)=\bigl\|j_{q,p}[x]_q\bigr\|
 \le\sum_j|a_j|\,|\ell_j([x]_q)|
 \le h(x).
\]

Weźmy teraz dowolną ciągłą seminormę hilbertowską $h$.
Jeżeli $h=0$, wybieramy $r=0$ i żądany operator jest zerowy. W przeciwnym razie
wybieramy $q\ge h$ i nuklearne przejście $j_{q,h}:E_q\to H_h$.
Niech $j_{q,h}z=\sum_ja_j\ell_j(z)y_j$ będzie jego
znormalizowanym przedstawieniem, z $A=\sum_j|a_j|>0$.
Definiujemy $r(x)^2=A\sum_j|a_j|\,|\ell_j([x]_q)|^2$.
Powyższy argument daje $r\ge h$ i ciągłość $r$.
Na $E/\ker r$ określamy
\[
 U([x]_r)=(\sqrt{A|a_j|}\,\ell_j([x]_q))_j\in\ell^2,
\qquad
 Ve_j=A^{-1/2}\sqrt{|a_j|}\,\operatorname{sgn}(a_j)y_j,
\]
gdzie $\operatorname{sgn}(a_j)=a_j/|a_j|$ dla $a_j\ne0$,
a $\operatorname{sgn}(0)=0$. Pierwsze odwzorowanie jest dobrze określone
na klasach, ponieważ dla $x\in\ker r$ wszystkie
ważone współrzędne są zerowe. Jest izometrią
i przedłuża się do $U:H_r\to\ell^2$.
Drugie odwzorowanie, początkowo określone na ciągach
o skończonym nośniku, spełnia
$\sum_j\|Ve_j\|^2\le A^{-1}\sum_j|a_j|=1$.
Dla takiego ciągu $z=(z_j)$ mamy
\[
 \Big\|\sum_j z_jVe_j\Big\|
 \le\|z\|_2\Big(\sum_j\|Ve_j\|^2\Big)^{1/2}.
\]
Zatem $V$ przedłuża się do ograniczonego operatora
$\ell^2\to H_h$, który jest operatorem
Hilberta--Schmidta na mocy powyższej sumy.
Na gęstym obrazie $E$ zachodzi
\[
 VU([x]_r)=\sum_j a_j\ell_j([x]_q)y_j=[x]_h.
\]
Wobec ciągłości $VU=j_{r,h}$.
Obraz izometrii $U$ jest domknięty, a obraz dowolnej
bazy ortonormalnej $H_r$ przez $U$ można uzupełnić
do bazy ortonormalnej $\ell^2$. Stąd
$\|VU\|_{\mathrm{HS}}^2\le\|V\|_{\mathrm{HS}}^2$,
co dowodzi żądanej własności.

\emph{(ii)$\Rightarrow$(i).} Dla dowolnej ciągłej seminormy
$p$ wybieramy $h\ge p$ hilbertowską, a następnie
hilbertowskie $r\ge h$ oraz $s\ge r$ tak, aby
$H_s\to H_r$ i $H_r\to H_h$ były operatorami
Hilberta--Schmidta. Ich złożenie $H_s\to H_h$ jest
nuklearne na mocy Lematu~\ref{lem:hs-compose-nuclear23}.
Nierówność $p\le h$ określa ograniczone odwzorowanie
$H_h\to E_p$. Na mocy Stwierdzenia~\ref{prop:nuclear-ideal23}
złożenie operatora nuklearnego z ograniczonym
jest nuklearne, zatem $j_{s,p}:E_s=H_s\to E_p$ jest
nuklearny (identyfikujemy złożenie na gęstym obrazie $E$). To jest warunek z Definicji~\ref{def:nuclear-lcs23}.
\end{proof}

\begin{twierdzenie}[Trwałość względem podprzestrzeni i ilorazów]
\label{thm:nuclear-subquot23}
Niech $E$ będzie nuklearną przestrzenią lokalnie wypukłą
Hausdorffa. Każda podprzestrzeń liniowa $F\subseteq E$
z topologią indukowaną jest nuklearna. Jeśli $M\subseteq E$
jest domkniętą podprzestrzenią liniową, to iloraz $E/M$
z topologią ilorazową także jest nuklearny.
\end{twierdzenie}
\begin{proof}
Wykorzystamy dwie własności operatorów Hilberta--Schmidta.
Jeżeli $J:H_1\to H_0$ jest operatorem Hilberta--Schmidta,
to jego ograniczenie do domkniętej podprzestrzeni $R\subseteq H_1$
z wartościami w $\overline{J(R)}$ też ma tę własność:
wystarczy dopełnić bazę ortonormalną $R$ do bazy $H_1$.
Jeżeli natomiast $N\subseteq H_1$ i $L\subseteq H_0$
są domkniętymi podprzestrzeniami oraz $J(N)\subseteq L$,
to operator na ilorazach $H_1/N\to H_0/L$ jest
Hilberta--Schmidta. Istotnie, po identyfikacjach ilorazów
z dopełnieniami ortogonalnymi jest to
$P_{L^\perp}J|_{N^\perp}$.
Dla bazy ortonormalnej $(u_i)$ podprzestrzeni $N^\perp$
mamy $\sum_i\|P_{L^\perp}Ju_i\|^2\le\sum_i\|Ju_i\|^2
\le\|J\|_{\mathrm{HS}}^2$, uzupełniając bazę do bazy $H_1$.
Nie potrzeba założenia $J(N)=L$.

Niech $F\subseteq E$ będzie dowolną podprzestrzenią.
Dla każdej ciągłej seminormy $p_F$ na $F$ istnieje
ciągła seminorma $p_E$ na $E$ taka, że
$p_F\le p_E|_F$. Aby to wykazać, wybieramy otwarte, absolutnie
wypukłe otoczenie zera $U$ w $E$ z
$U\cap F\subseteq\{x\in F:p_F(x)<1\}$ i bierzemy
jego ciągły funkcjonał Minkowskiego $p_E$.
Dla $x\in F$ i $t>p_E(x)$ mamy $x/t\in U\cap F$,
więc $p_F(x)<t$. Przechodząc do infimum po $t$,
otrzymujemy $p_F(x)\le p_E(x)$. Z kryterium hilbertowskiego
(Twierdzenie~\ref{thm:hilbert-nuclear23}) wybieramy
hilbertowskie $h\ge p_E$, $r\ge h$, $s\ge r$,
dla których $H_s\to H_r$ i $H_r\to H_h$ są
Hilberta--Schmidta. Dla $t\in\{h,r,s\}$ przestrzeń
$H_{t|_F}$ identyfikuje się izometrycznie z domknięciem
obrazu $F$ w $H_t$: odwzorowanie $[x]_{t|_F}\mapsto[x]_t$
jest izometrią na gęstej podprzestrzeni tego domknięcia.
Przejścia kanoniczne przenoszą te domknięcia do
odpowiednich domknięć na mocy ciągłości i zgodności na $F$. Wskazana własność ograniczeń zapewnia,
że $H_{s|_F}\to H_{r|_F}$ oraz
$H_{r|_F}\to H_{h|_F}$ są Hilberta--Schmidta.
Ich złożenie jest nuklearne. Ponieważ
$p_F\le h|_F$, istnieje ograniczone odwzorowanie
$H_{h|_F}\to F_{p_F}$; jego złożenie z poprzednim
przejściem dowodzi nuklearności $F$.

Niech teraz $M\subseteq E$ będzie domknięta i
$Q=E/M$ z odwzorowaniem ilorazowym $\pi:E\to Q$.
Dla ciągłej seminormy $p_Q$ na $Q$ zastosujmy kryterium
do $p_Q\circ\pi$ i wybierzmy hilbertowskie seminormy
$h\ge p_Q\circ\pi$, $r\ge h$, $s\ge r$,
których kolejne przejścia są Hilberta--Schmidta.
Dla $t\in\{h,r,s\}$ kładziemy
\[
 \bar t(\pi x)=\inf_{m\in M}t(x+m).
\]
Wartość nie zależy od reprezentanta klasy $\pi x$.
Niech $N_t$ oznacza domknięcie obrazu $M$ w $H_t$.
Wówczas
\[
 \bar t(\pi x)=\operatorname{dist}([x]_t,N_t)
 =\|P_{N_t^\perp}[x]_t\|.
\]
Ostatnia równość wynika z twierdzenia o rzucie
ortogonalnym. Jest to seminorma hilbertowska, zadana
przez iloczyn skalarny obrazów przez
$\pi x\mapsto P_{N_t^\perp}[x]_t$.
Jej ciągłość wynika z równości
$\{y:\bar t(y)<\varepsilon\}=
\pi(\{x:t(x)<\varepsilon\})$ i otwartości $\pi$.
Odwzorowanie $Q/\ker\bar t\to H_t/N_t$ dane przez
$[\pi x]_{\bar t}\mapsto [x]_t+N_t$ jest izometryczne
i ma gęsty obraz, ponieważ obraz $E$ w $H_t$ jest gęsty.
Przedłuża się więc do izometrycznego izomorfizmu
$H_{\bar t}\simeq H_t/N_t$.
Ponadto $\bar s\ge\bar r\ge\bar h\ge p_Q$,
gdyż $h(x+m)\ge p_Q(\pi(x+m))=p_Q(\pi x)$
dla każdego $m\in M$.
Operator $H_s\to H_r$ odwzorowuje $N_s$ w $N_r$,
podobnie $H_r\to H_h$ odwzorowuje $N_r$ w $N_h$.
Z własności operatorów na ilorazach otrzymujemy
przejścia Hilberta--Schmidta
$H_{\bar s}\to H_{\bar r}\to H_{\bar h}$.
Ich złożenie jest nuklearne, zaś
$H_{\bar h}\to Q_{p_Q}$ jest ograniczone. Zatem
$j_{\bar s,p_Q}:Q_{\bar s}\to Q_{p_Q}$ jest nuklearny,
co kończy dowód.
\end{proof}

\begin{lemat}[Nuklearność przeliczalnej sumy prostej]
\label{lem:nuclear-direct-sum23}
Niech $E_1,E_2,\ldots$ będą nuklearnymi przestrzeniami
lokalnie wypukłymi Hausdorffa. Wówczas ich lokalnie wypukła
suma prosta $E=\bigoplus_{n\ge1}E_n$, z najsilniejszą
topologią lokalnie wypukłą, dla której wszystkie kanoniczne
włożenia $E_n\to E$ są ciągłe, jest nuklearna.
\end{lemat}
\begin{proof}
Suma prosta jest Hausdorffa: dla niezerowego
$x=(x_n)$ wybieramy niezerową współrzędną $x_n$
i ciągłą seminormę na $E_n$ niezerową na $x_n$.
Jej złożenie z rzutem na $n$-tą współrzędną jest
ciągłe na $E$, ponieważ jego ograniczenia do wszystkich
składników są ciągłe.
Niech $p$ będzie ciągłą seminormą na $E$,
$p_n=p|_{E_n}$. Dla każdego $n$ z nuklearności $E_n$
wybieramy ciągłą seminormę $q_n^0\ge p_n$, dla której
$j_{q_n^0,p_n}$ jest nuklearny. Zastępując $q_n^0$
przez $q_n=c_nq_n^0$ dla dostatecznie dużego $c_n\ge1$,
możemy zapewnić
$\|j_{q_n,p_n}\|_{\mathcal N}<2^{-n}$:
przeskalowanie normy w przestrzeni źródłowej przez
$c_n$ dzieli normę każdego funkcjonału przez $c_n$,
więc dzieli także normę nuklearną przejścia przez $c_n$
(obie nierówności wynikają ze skalowania w obie strony).
Definiujemy na algebraicznej sumie prostej
$q((x_n))=\sum_nq_n(x_n)$; jest to suma skończona.
Seminorma $q$ jest ciągła, ponieważ jej ograniczenie
do każdego $E_n$ jest ciągłe, a topologia sumy prostej
jest końcową topologią lokalnie wypukłą. Z nierówności
trójkąta wynika $p\le\sum_n p_n\le q$.

Mamy $\ker q=\bigoplus_n\ker q_n$, a norma na ilorazie
$E/\ker q$ jest sumą norm współrzędnych.
Uzupełnienie $E/\ker q$ identyfikuje się izometrycznie
z $\bigoplus_n^{\ell^1}(E_n)_{q_n}$: jest to przestrzeń
Banacha, a obcięcie ogona i przybliżenie skończonej
liczby współrzędnych elementami $E_n/\ker q_n$
dowodzi gęstości obrazu $E/\ker q$.
Dla ustalonego $n$ kanoniczny operator
$(E_n)_{q_n}\to E_p$ jest złożeniem nuklearnego
$j_{q_n,p_n}$ i izometrycznego osadzenia
$(E_n)_{p_n}\to E_p$: na gęstym ilorazie
norma obrazu jest dokładnie $p_n(x_n)$.
Ze Stwierdzenia~\ref{prop:nuclear-ideal23} (własność
ideałowa) wynika, że norma nuklearna złożenia jest
mniejsza niż $2^{-n}$.
Wybieramy jego przedstawienie nuklearne, którego suma
norm składników jest mniejsza niż $2^{-n}$.
Każdy funkcjonał występujący w tym przedstawieniu
przedłużamy na sumę $\ell^1$ przez złożenie
z rzutem współrzędnym normy co najwyżej $1$ (rzut
na zerowy składnik ma normę $0$).
Indeksy par $(n,j)$ tworzą zbiór przeliczalny.
Po połączeniu wszystkich przedstawień w jeden szereg
suma iloczynów norm funkcjonałów i wektorów jest
ograniczona przez $\sum_{n\ge1}2^{-n}<\infty$.
Szereg wyznacza więc operator nuklearny bez potrzeby
odwoływania się do zupełności przestrzeni operatorów
nuklearnych. Ich suma jest dokładnie
kanonicznym operatorem $E_q\to E_p$:
najpierw sprawdzamy tę równość dla wektorów o skończonym
nośniku, a następnie korzystamy z ciągłości i gęstości.
To dowodzi nuklearności $E$.
\end{proof}

\begin{twierdzenie}[Trwałość względem przeliczalnych granic induktywnych]
\label{thm:nuclear-indlim23}
Niech $(E_n,u_n)_{n\ge1}$ będzie ciągiem nuklearnych
przestrzeni lokalnie wypukłych Hausdorffa i ciągłych
odwzorowań liniowych $u_n:E_n\to E_{n+1}$.
Jeśli lokalnie wypukła granica induktywna
$E=\varinjlim_n E_n$ jest Hausdorffa, to $E$ jest nuklearna.
Nie zakładamy metryzowalności $E$ ani ścisłości granicy.
\end{twierdzenie}
\begin{proof}
Niech $S=\bigoplus_{n\ge1}E_n$ będzie lokalnie wypukłą
sumą prostą, a $\iota_n:E_n\to E$ odwzorowaniami
kanonicznymi. Definiujemy
$\Pi:S\to E$ przez
$\Pi((x_n))=\sum_n\iota_n(x_n)$; suma jest skończona.
Z końcowej własności topologii sumy prostej wynika
ciągłość $\Pi$, a z definicji granicy induktywnej
wynika jej surjektywność.

Topologia ilorazowa na $S/\ker\Pi$ ma dokładnie tę samą
własność uniwersalną co lokalnie wypukła topologia
granicy induktywnej: odwzorowanie liniowe $T:E\to G$
do dowolnej przestrzeni lokalnie wypukłej $G$ jest ciągłe
wtedy i tylko wtedy, gdy wszystkie $T\circ\iota_n$ są ciągłe.
Zatem $E\simeq S/\ker\Pi$ jako przestrzenie lokalnie
wypukłe. Ponieważ $E$ jest Hausdorffa,
$\ker\Pi=\Pi^{-1}(\{0\})$ jest domknięte w $S$.
Suma prosta $S$ jest nuklearna na mocy
Lematu~\ref{lem:nuclear-direct-sum23}, a iloraz
przez domkniętą podprzestrzeń jest nuklearny na mocy
Twierdzenia~\ref{thm:nuclear-subquot23}. Wynika stąd
nuklearność $E$. Ścisłość i metryzowalność granicy
nie zostały wykorzystane.
\end{proof}

\begin{twierdzenie}[Nuklearność $\mathcal D_K$ i $\mathcal D$]
\label{thm:nuclear-D23}
Niech $d\ge1$, $\Omega\subseteq\R^d$ będzie niepustym
zbiorem otwartym, a $K\Subset\Omega$ zbiorem zwartym.
Wówczas $\mathcal D_K(\Omega)$ ze standardową topologią
Frécheta i $\mathcal D(\Omega)$ ze standardową topologią
LF są nuklearnymi przestrzeniami lokalnie wypukłymi.
\end{twierdzenie}
\begin{proof}
Dla $K=\varnothing$ przestrzeń $\mathcal D_K(\Omega)$
jest zerowa, a więc nuklearna. Dla niepustego $K$
rozważmy najpierw $\mathcal D_K(\Omega)$ jako
podprzestrzeń $\mathcal E(\Omega)=C^\infty(\Omega)$.
Ponieważ $K$ jest domknięty w $\Omega$,
warunek $\supp f\subseteq K$ jest równoważny
$f(x)=0$ dla wszystkich $x\in\Omega\setminus K$.
A zatem
\[
 \mathcal D_K(\Omega)=
 \bigcap_{x\in\Omega\setminus K}\ker(\mathrm{ev}_x),
 \qquad \mathrm{ev}_x(f)=f(x).
\]
Każda ewaluacja jest ciągła w topologii $C^\infty$,
więc $\mathcal D_K(\Omega)$ jest podprzestrzenią
domkniętą. Ponadto wszystkie pochodne jej elementów
znikają poza $K$, stąd dla dowolnego $L\Subset\Omega$
\[
 p_{L,m}(f)\le p_{K,m}(f),\qquad
 f\in\mathcal D_K(\Omega).
\]
Ponieważ $K$ samo jest zwarte, rodzina $p_{K,m}$
jest zarazem rodziną ograniczeń seminorm z $C^\infty(\Omega)$.
Topologia indukowana jest więc dokładnie standardową
topologią Frécheta na $\mathcal D_K(\Omega)$.
Twierdzenia~\ref{thm:nuclear-cinfty23} i
\ref{thm:nuclear-subquot23} dowodzą jej nuklearności.

Wybierzmy wyczerpanie $\Omega$ zbiorami zwartymi
$K_1\subseteq\operatorname{int}K_2\subseteq\cdots\Subset\Omega$
takie, że każdy zwarty podzbiór $\Omega$ zawiera się
w pewnym $K_n$. Definicja~\ref{def:D-LF18} identyfikuje standardową
topologię LF na $\mathcal D(\Omega)$ z topologią
lokalnie wypukłej granicy induktywnej
$\varinjlim_n\mathcal D_{K_n}(\Omega)$, a
Twierdzenie~\ref{thm:D-stage18} zapewnia warunek Hausdorffa.
Włożenia kolejnych etapów są ciągłe.
Jej nuklearność wynika teraz z
Twierdzenia~\ref{thm:nuclear-indlim23}.
\end{proof}

\begin{uwaga}[Dlaczego sam argument z $C^\infty$ nie wystarcza]
\label{rem:nuclear-D-summary23}
Nuklearność $\mathcal D_K$ jest konsekwencją trwałości
względem podprzestrzeni, ale nuklearność
$\mathcal D=\varinjlim\mathcal D_{K_n}$ wymaga osobnego
wyniku dotyczącego granic induktywnych. Powyższe dowody
ujawniają wszystkie użyte przejścia. Nie wynika z tego
równość topologii induktywnej tensorowej $i$ z topologią
rzutową $\pi$: dotyczą one innych własności uniwersalnych.
\end{uwaga}

\subsection{Twierdzenie o jądrze Schwartza}

\noindent\emph{Powiązania tematyczne w Bibliotece matematycznej:} \atlasref{LA-10::thm:lin10-istnienie}{LA-10: algebraiczna konstrukcja iloczynu tensorowego (przed uzupełnieniem topologicznym)}.


\begin{definicja}[Forma dwuliniowa zadana jądrem]
\label{def:kernel-pair23}
Niech $\Omega\subseteq\R^{d_1}$ i
$\Lambda\subseteq\R^{d_2}$ będą niepustymi zbiorami otwartymi.
Dla $K\in\mathcal D'(\Omega\times\Lambda)$ definiujemy
formę na parach funkcji testowych przez
$b_K(f,g)=\langle K,f\otimes g\rangle$,
dla $f\in\mathcal D(\Omega)$ i $g\in\mathcal D(\Lambda)$,
gdzie $(f\otimes g)(x,y)=f(x)g(y)$.
Parowanie jest dwuliniowe także nad $\C$; nie stosujemy
sprzężenia zespolonego.
\end{definicja}

\begin{lemat}[Rozkład testu na szeregi tensorowe]
\label{lem:tensor-density23}
Niech $d_1,d_2\ge1$, a $\Omega\subseteq\R^{d_1}$ i
$\Lambda\subseteq\R^{d_2}$ będą niepustymi zbiorami otwartymi.
Dla każdego $\Phi\in\mathcal D(\Omega\times\Lambda)$
istnieje przedstawienie $\Phi$ jako sumy skończenie
wielu szeregów postaci
$\sum_{n\in\mathbb Z^{d_2}}f_n\otimes g_n$,
zbieżnych w $\mathcal D(\Omega\times\Lambda)$,
gdzie wszystkie $f_n$ mają nośniki w jednym zwartym
$K\Subset\Omega$, wszystkie $g_n$ w jednym zwartym
$L\Subset\Lambda$ i dla każdego $m,r\in\N_0$
\[
 \sum_n p_{K,m}(f_n)\,p_{L,r}(g_n)<\infty.
\]
W szczególności algebraiczna przestrzeń tensorowa
$\mathcal D(\Omega)\otimes\mathcal D(\Lambda)$
jest gęsta w $\mathcal D(\Omega\times\Lambda)$.
\end{lemat}
\begin{proof}
Wybieramy lokalnie skończone gładkie rozkłady jedności
$(\eta_i)$ na $\Omega$ i $(\theta_j)$ na $\Lambda$
podporządkowane pokryciom stosunkowo zwartymi otwartymi
sześcianami. Dla każdego $i,j$ ustalamy zbiory zwarte
$K_i=\supp\eta_i$ i $L_j=\supp\theta_j$ oraz otwarty
sześcian $V_j$ z $L_j\subset V_j\Subset\Lambda$.
Wybieramy rzeczywistą $\chi_j\in\mathcal D(V_j)$,
równą $1$ w otoczeniu $L_j$ (Lemat~\ref{lem:cutoffs17}).
Funkcję $\Phi_{ij}=\eta_i\theta_j\Phi$ przedłużamy
zerem w zmiennej $y$ poza $V_j$ i okresowo z jego
domknięcia $Q_j$; przedłużenie jest gładkie,
gdyż nośnik w $y$ leży w $L_j$.
Niech $a_j$ będzie środkiem $Q_j$, a $2R_j$ jego
długością boku. Kładziemy
\[
 e_{jn}(y)=e^{i\pi n\cdot(y-a_j)/R_j},
 \quad
 f_{ijn}(x)=\frac{1}{(2R_j)^{d_2}}
 \int_{Q_j}\Phi_{ij}(x,y)\overline{e_{jn}(y)}
 \,\mathrm d\lambda_{d_2}(y),
 \quad g_{jn}=\chi_j e_{jn}.
\]
Funkcje $f_{ijn}$ są gładkie i mają nośniki w $K_i$:
pochodne względem $x$ obliczamy pod całką po zwartym
sześcianie. Niech $L_j^*=\supp\chi_j\cup L_j$.
Dla $m,r,t\in\N_0$, całkując przez części za pomocą
$(1-\Delta_y)^t$, otrzymujemy
\[
 \begin{aligned}
 p_{K_i,m}(f_{ijn})&\le C_{j,t}
 (1+|n|)^{-2t}p_{K_i\times L_j^*,m+2t}(\Phi_{ij}),\\
 p_{L_j^*,r}(g_{jn})&\le C_{j,r}'(1+|n|)^r .
 \end{aligned}
\]
Pierwsza nierówność wynika z tożsamości
$(1-\Delta_y)e_{jn}=(1+\pi^2|n|^2/R_j^2)e_{jn}$,
braku wyrazów brzegowych oraz ograniczenia pochodnych
funkcji podcałkowej; druga ze wzoru Leibniza.
Przy $2t>r+d_2$ szereg iloczynów seminorm jest zbieżny
(liczba punktów kratowych o normie maksimum $s$
jest rzędu co najwyżej $(1+s)^{d_2-1}$).
W szczególności dla każdego $q$ szereg
$\sum_n f_{ijn}\otimes g_{jn}$ jest bezwzględnie
zbieżny w seminormie $p_{K_i\times L_j^*,q}$,
stosując powyższe oszacowanie z $m=r=q$.
Zupełność etapu $\mathcal D_{K_i\times L_j^*}$
(Twierdzenie~\ref{thm:DK-frechet18}) daje granicę.

Dla ustalonego $x$ szereg $\sum_n f_{ijn}(x)e_{jn}(y)$
na $Q_j$ jest jednostajnie zbieżny, a jego granicę
identyfikujemy z $\Phi_{ij}(x,\cdot)$ dzięki
zupełności układu trygonometrycznego w $L^2(Q_j)$
i ciągłości obu funkcji. Po pomnożeniu przez $\chi_j$
równość obowiązuje na całej $\Lambda$: na $Q_j$
mamy $\chi_j\Phi_{ij}=\Phi_{ij}$, a poza $Q_j$
obie strony znikają. Warunek $\supp\chi_j\subset V_j$
jest istotny: eliminuje inne kopie okresowe.
Stąd $\Phi_{ij}=\sum_n f_{ijn}\otimes g_{jn}$.
Tylko skończenie wiele par $(i,j)$ daje niezerowe
$\Phi_{ij}$, dzięki lokalnej skończoności i zwartości
projekcji nośnika $\Phi$. Sumujemy te rozwinięcia,
biorąc sumy mnogościowe odpowiednich zwartych zbiorów
$K_i$ i $L_j^*$. Daje to wspólny zwarty nośnik
i żądaną gęstość.

Nad $\R$ w powyższych szeregach dla rzeczywistej $\Phi$
zastępujemy każdy składnik przez jego część rzeczywistą:
$\operatorname{Re}f_{ijn}\otimes\operatorname{Re}g_{jn}
-\operatorname{Im}f_{ijn}\otimes\operatorname{Im}g_{jn}$.
Suma pozostaje zbieżna z tymi samymi oszacowaniami
(z dodatkowym czynnikiem $2$).
\end{proof}

\begin{lemat}[Ciągłość formy na ustalonych etapach]
\label{lem:separate-joint23}
Niech $E,F$ będą przestrzeniami Frécheta nad
$\K\in\{\R,\C\}$, a $b:E\times F\to\K$
formą dwuliniową ciągłą osobno w każdej zmiennej.
Wtedy $b$ jest łącznie ciągła.
\end{lemat}
\begin{proof}
Wybierzmy malejącą przeliczalną bazę
$(V_n)$ otoczeń zera w $F$, złożoną ze zbiorów
zrównoważonych. Niech
$A_n=\{x\in E: |b(x,y)|\le1\ \text{dla wszystkich }y\in V_n\}$.
Każdy $A_n$ jest domknięty, ponieważ $x\mapsto b(x,y)$
jest ciągły dla ustalonego $y$. Z ciągłości w drugiej
zmiennej wynika $E=\bigcup_n A_n$.
Z twierdzenia Baire'a pewien $A_n$ ma niepuste wnętrze;
ponieważ $A_n$ jest wypukły, zrównoważony i domknięty,
ma wówczas otoczenie zera we wnętrzu. Istotnie,
jeśli $x_0+U\subseteq A_n$ dla otwartego symetrycznego
otoczenia zera $U$, to również $-x_0+U\subseteq A_n$;
wypukłość daje $U\subseteq A_n$.
Stąd $|b(x,y)|\le1$ na $U\times V_n$.
Skalowanie daje ciągłość w $(0,0)$, a dla
dowolnego $(x_0,y_0)$ korzystamy z rozwinięcia
$b(x_0+x,y_0+y)-b(x_0,y_0)
=b(x_0,y)+b(x,y_0)+b(x,y)$ i ciągłości osobnej.
\end{proof}

\begin{uwaga}[Ciągłość osobna i łączna]
\label{rem:kernel-separate23}
Dla formy $b:\mathcal D(\Omega)\times\mathcal D(\Lambda)\to\K$
ciągłość osobna nie musi implikować łącznej ciągłości na
całym iloczynie przestrzeni LF. Implikacja ta jest jednak
prawdziwa po ograniczeniu formy do
$\mathcal D_K(\Omega)\times\mathcal D_L(\Lambda)$,
gdyż obie przestrzenie są przestrzeniami Frécheta.
To rozróżnienie jest istotne w twierdzeniu o jądrze;
właściwą globalną topologią tensorową jest topologia
induktywna, omówiona w następnym rozdziale.
\end{uwaga}

\begin{twierdzenie}[Twierdzenie o jądrze Schwartza]
\label{thm:schwartz-kernel23}
Niech $\Omega\subseteq\R^{d_1}$ i
$\Lambda\subseteq\R^{d_2}$ będą niepustymi zbiorami
otwartymi. Dla każdego ciągłego operatora liniowego
\[
 A:\mathcal D(\Omega)\longrightarrow
 \mathcal D'(\Lambda)_\beta
\]
istnieje dokładnie jedna dystrybucja
$K_A\in\mathcal D'(\Omega\times\Lambda)$ taka, że
\[
 \langle Af,g\rangle=\langle K_A,f\otimes g\rangle
 \qquad(f\in\mathcal D(\Omega),\ g\in\mathcal D(\Lambda)).
\]
Odwrotnie, każda dystrybucja
$K\in\mathcal D'(\Omega\times\Lambda)$ określa w ten
sposób ciągły operator liniowy
$A_K:\mathcal D(\Omega)\to\mathcal D'(\Lambda)_\beta$.
Indeks $\beta$ oznacza silną topologię przestrzeni dualnej,
czyli jednostajną zbieżność na ograniczonych zbiorach
funkcji testowych.
\end{twierdzenie}
\begin{proof}
Niech $b(f,g)=\langle Af,g\rangle$. Ciągłość $A$
i ciągłość ewaluacji na pojedynczej funkcji testowej
zapewniają osobną ciągłość $b$. Na przestrzeniach
$\mathcal D_{K_x}(\Omega)$ i $\mathcal D_{K_y}(\Lambda)$
forma $b$ jest wspólnie ciągła na mocy
Lematu~\ref{lem:separate-joint23}; zatem istnieją
$C,m,r$ takie, że
$|b(f,g)|\le C p_{K_x,m}(f)p_{K_y,r}(g)$.
Zbudujemy jeden liniowy funkcjonał na całej przestrzeni
testowej, a następnie sprawdzimy jego niezależność
od wyborów pomocniczych. Wybierzmy lokalnie skończone
rozkłady jedności $(\eta_i)$ na $\Omega$ oraz
$(\theta_j)$ na $\Lambda$, których elementy mają zwarte
nośniki w sześcianach stosunkowo zwartych. Produkty
$\eta_i(x)\theta_j(y)$ tworzą rozkład jedności na
$\Omega\times\Lambda$. Dla danej
$\Phi\in\mathcal D(\Omega\times\Lambda)$ tylko skończenie
wiele funkcji $\Phi_{ij}=\eta_i\theta_j\Phi$ jest niezerowych.
Dla każdej pary $(i,j)$ wybieramy z góry odcięcie
$\chi_j$ i układ Fouriera z dowodu
Lematu~\ref{lem:tensor-density23}, niezależnie od $\Phi$.
Otrzymujemy $\Phi_{ij}=\sum_n f_{ijn}\otimes g_{jn}$,
gdzie $f_{ijn}$ zależy liniowo od $\Phi$,
a $g_{jn}$ nie zależy od $\Phi$.
Nad $\R$ stosujemy opisane w lemacie dwa rzeczywiste
składniki; dalej zapisujemy wzory w przypadku zespolonym.
Definiujemy
\[
 \langle K_A,\Phi\rangle=\sum_{i,j}\sum_n b(f_{ijn},g_{jn}),
\]
gdzie suma po $i,j$ jest skończona, a wszystkie szeregi
wewnętrzne są bezwzględnie zbieżne na mocy oszacowań
seminormowych i ciągłości $b$ na odpowiednich etapach.
Konstrukcja jest liniowa, ponieważ wybory odcięć i baz
Fouriera nie zależą od $\Phi$, a współczynniki Fouriera
zależą od niej liniowo.

Aby sprawdzić ciągłość, ustalmy $M\Subset\Omega\times\Lambda$.
Lokalna skończoność rozkładów jedności daje skończoną
liczbę par $(i,j)$, które mogą wystąpić dla
$\Phi\in\mathcal D_M$. Dla każdej takiej pary istnieją $C_{ij},m_{ij},r_{ij}$
z oszacowania dwuliniowego na $\mathcal D_{K_i}
\times\mathcal D_{L_j^*}$. Wybieramy $t_{ij}$ z
$2t_{ij}>r_{ij}+d_2$. Z oszacowań w dowodzie
Lematu~\ref{lem:tensor-density23} i wzoru Leibniza
dla $\eta_i\theta_j$ otrzymujemy
\[
 |b(f_{ijn},g_{jn})|\le C'_{ij}
 (1+|n|)^{r_{ij}-2t_{ij}}
 p_{M,m_{ij}+2t_{ij}}(\Phi).
\]
Pochodne $\Phi$ znikają poza $M$, zatem po prawej
wystarcza seminorma na $M$. Sumujemy po $n$,
a następnie po skończonej liczbie par i bierzemy
$r_M=\max_{i,j}(m_{ij}+2t_{ij})$ (dla pustej
rodziny par funkcjonał jest zerowy). Otrzymujemy
stałą $C_M>0$ i $r_M\in\N_0$ takie, że
\[
 |\langle K_A,\Phi\rangle|\le C_M p_{M,r_M}(\Phi)
 \qquad (\Phi\in\mathcal D_M(\Omega\times\Lambda)).
\]
Jest to ciągłość na każdym etapie, więc własność
uniwersalna granicy induktywnej daje
$K_A\in\mathcal D'(\Omega\times\Lambda)$.

Dla $\Phi=f\otimes g$ współczynniki mają szczególną
postać $f_{ijn}=(\eta_i f)c_{jn}(\theta_j g)$,
gdzie $c_{jn}$ jest skalarnym współczynnikiem Fouriera
w $Q_j$. Z dowodu lematu
$\sum_n c_{jn}(\theta_j g)g_{jn}=\theta_j g$
w ustalonym etapie $\mathcal D_{L_j^*}$.
Liniowość i ciągłość $b$ względem drugiego argumentu
przy ustalonym $\eta_i f$ dają
\[
 \sum_n b(f_{ijn},g_{jn})=b(\eta_i f,\theta_j g).
\]
Nad $\R$ tę samą równość otrzymujemy, rozwijając
rzeczywistą część szeregu i korzystając z rzeczywistej
dwuliniowości. Sumując skończenie wiele składników
rozkładów jedności, dostajemy
$\langle K_A,f\otimes g\rangle=b(f,g)$.
Jeżeli zmienimy rozkłady jedności i odcięcia, otrzymamy
inny ciągły funkcjonał na całej przestrzeni $\mathcal D$,
ale o tych samych wartościach na tensorach elementarnych.
Ich liniowa powłoka jest gęsta na mocy
Lematu~\ref{lem:tensor-density23}; dwa funkcjonały muszą
więc być równe. To zarazem uzasadnia jednoznaczność,
bez zakładania gęstości tensorów o nośnikach w arbitralnym
zwartym zbiorze $M$.

W odwrotną stronę, dla $K\in\mathcal D'(\Omega\times\Lambda)$
forma $b_K(f,g)=\langle K,f\otimes g\rangle$ jest
łącznie ciągła na każdym
$\mathcal D_{K_x}(\Omega)\times\mathcal D_{K_y}(\Lambda)$,
ponieważ dla ustalonych niepustych zwartych $K_x,K_y$
z ciągłości $K$ na etapie $\mathcal D_{K_x\times K_y}$
otrzymujemy pewne $C>0$ i $q\in\N_0$ z
\[
 |\langle K,f\otimes g\rangle|
 \le C p_{K_x\times K_y,q}(f\otimes g)
 \le C p_{K_x,q}(f)p_{K_y,q}(g).
\]
Ostatnia nierówność wynika z rozdzielenia pochodnych
względem $x$ i $y$. Dla ustalonego $f$ wybieramy
$K_x$ zawierający jego nośnik. Forma $g\mapsto b_K(f,g)$
jest liniowa i ciągła na każdym etapie $\mathcal D_{K_y}$,
więc na mocy Twierdzenia~\ref{thm:LF-universal18}
wyznacza $A_Kf\in\mathcal D'(\Lambda)$. Dwuliniowość
zapewnia liniowość $A_K$.
Aby sprawdzić ciągłość w topologii silnej, niech
$B\subseteq\mathcal D(\Lambda)$ będzie ograniczony.
Z Twierdzenia~\ref{thm:D-bounded18} istnieje
$K_y\Subset\Lambda$, dla którego $B$ jest ograniczony
w $\mathcal D_{K_y}$. Na każdym etapie $\mathcal D_{K_x}$
mamy, dzięki wspólnej ciągłości formy,
\[
 \sup_{g\in B}|\langle A_Kf,g\rangle|
 \le C\,p_{K_x,m}(f)\sup_{g\in B}p_{K_y,r}(g).
\]
To dowodzi ciągłości $A_K$ na każdym etapie, a więc,
przez własność uniwersalną granicy induktywnej,
na całej $\mathcal D(\Omega)$. Jednoznaczność $K_A$
wynika ponownie z gęstości tensorów elementarnych.
\end{proof}

\begin{przyklad}[Jądro operatora o gładkim jądrze]
\label{ex:smooth-kernel23}
Jeżeli $k\in C^\infty(\Omega\times\Lambda)$, to dystrybucja
regularna $T_k$ jest jądrem operatora
$(Af)(y)=\int_\Omega k(x,y)f(x)\,\mathrm d\lambda_{d_1}(x)$.
Dla $g\in\mathcal D(\Lambda)$ mamy
\[
\langle Af,g\rangle
=\int_{\Lambda}\int_{\Omega}k(x,y)f(x)g(y)
\,\mathrm d\lambda_{d_1}(x)\,\mathrm d\lambda_{d_2}(y).
\]
Dla ustalonego $f$ różniczkowanie względem $y$
pod całką jest uzasadnione przez ograniczoność pochodnych
$k$ na iloczynie $\supp f$ i dowolnego zwartego
otoczenia punktu $y$; funkcja $Af$ jest gładka.
Twierdzenie Fubiniego stosujemy, ponieważ funkcja
podcałkowa jest ograniczona i znika poza
$\supp f\times\supp g$. Nie wymagamy, żeby $k$
miało zwarty nośnik. Identyfikując $Af$ z dystrybucją
regularną, z twierdzenia o jądrze otrzymujemy
ciągłość operatora $\mathcal D(\Omega)\to
\mathcal D'(\Lambda)_\beta$.
\end{przyklad}

\begin{przyklad}[Jądro osobliwe na przekątnej]
\label{ex:diagonal-kernel23}
Dla $\Omega=\Lambda\subseteq\R^d$ operator naturalnego
włożenia $\mathcal D(\Omega)\hookrightarrow\mathcal D'(\Omega)$
ma jądro $\delta_\Delta$ zadane przez
$\langle\delta_\Delta,\Phi\rangle
=\int_\Omega\Phi(x,x)\,\mathrm d\lambda_d(x)$.
Dla $\Phi\in\mathcal D_M(\Omega\times\Omega)$
całka jest ograniczona przez
$\lambda_d(\operatorname{pr}_1 M)p_{M,0}(\Phi)$,
co dowodzi ciągłości na każdym etapie.
Na tensorze $f\otimes g$ otrzymujemy $\int fg\,
\mathrm d\lambda_d$, czyli parowanie naturalnego
włożenia. Nośnik $\delta_\Delta$ jest zawarty
w przekątnej $\Delta=\{(x,x):x\in\Omega\}$;
każde otoczenie punktu $(x,x)$ zawiera iloczyn
małych kul o środku $x$, na którym można wybrać
nieujemny tensor $f\otimes f$ z $f\ne0$.
Wtedy $\langle\delta_\Delta,f\otimes f\rangle>0$,
więc nośnik jest równy $\Delta$.
Gdyby $\delta_\Delta=T_h$ dla
$h\in L^1_{\mathrm{loc}}(\Omega\times\Omega)$,
to ograniczenie do otwartego dopełnienia $\Delta$
byłoby zerowe. Różnowartościowość przyporządkowania
funkcji dystrybucji regularnej daje $h=0$ prawie
wszędzie poza $\Delta$. Przekątna ma miarę
$\lambda_{2d}$ zero, zatem $h=0$ prawie wszędzie,
co przeczy $\delta_\Delta\ne0$.
\end{przyklad}

\subsection{Uzupełnienie: funkcje Hermite'a i nuklearność przestrzeni Schwartza}
\label{sec:hermite-supp23}

\begin{lemat}[Rozwinięcie Hermite'a i seminormy]
\label{lem:hermite-schwartz23}
Niech $d\ge1$ będzie liczbą całkowitą. Dla $n\in\N_0$ definiujemy
\[
 P_n(x)=(-1)^n e^{x^2}\frac{\mathrm d^n}{\mathrm dx^n}e^{-x^2},
 \qquad h_n(x)=\frac{P_n(x)e^{-x^2/2}}{(2^n n!\sqrt\pi)^{1/2}}.
\]
Dla $\nu\in\N_0^d$ niech $h_\nu(x)=\prod_{j=1}^d h_{\nu_j}(x_j)$
oraz $|\nu|=\sum_j\nu_j$. Wówczas $(h_\nu)_\nu$
jest bazą ortonormalną $L^2(\R^d;\K)$, a odwzorowanie
\[
 H:\mathcal S(\R^d;\K)\to s(\N_0^d),\qquad
 H(f)=(\langle f,h_\nu\rangle_{L^2})_\nu
\]
jest izomorfizmem liniowo-topologicznym. Iloczyn skalarny
jest liniowy w pierwszym argumencie.
\end{lemat}
\begin{proof}
W jednym wymiarze z definicji $P_0=1$ i $P_{n+1}=2xP_n-P_n'$.
Indukcja daje stopień $n$ i współczynnik wiodący $2^n$.
Dla $n\ge m$, całkując $n$ razy przez części, otrzymujemy
\[
 \int_\R P_nP_m e^{-x^2}\,\mathrm d\lambda_1
 =\int_\R P_m^{(n)}e^{-x^2}\,\mathrm d\lambda_1.
\]
Wyrazy brzegowe znikają, gdyż pochodne Gaussa są
wielomianami pomnożonymi przez Gaussa. Wynik wynosi zero
dla $n>m$ oraz $2^n n!\sqrt\pi$ dla $n=m$.
Symetria daje ortogonalność również dla $n<m$.

Jeżeli $f\in L^2(\R;\C)$ jest ortogonalna do wszystkich
$h_n$, to $\int_\R f(x)x^n e^{-x^2/2}\,\mathrm d\lambda_1(x)=0$
dla każdego $n$, ponieważ $(P_0,\ldots,P_n)$ jest bazą
wielomianów stopnia co najwyżej $n$. Niech $g(x)=f(x)e^{-x^2/2}$.
Z nierówności Cauchy'ego--Schwarza mamy $g\in L^1\cap L^2$.
Dla każdego $R>0$ funkcja $e^{-x^2/2+R|x|}$ należy do $L^2$,
więc $|g(x)|e^{R|x|}$ jest całkowalna.
Rozwijając wykładniczą i stosując zbieżność zmajoryzowaną
dla $|z|\le R$, otrzymujemy
\[
 \int_\R g(x)e^{zx}\,\mathrm d\lambda_1(x)
 =\sum_{n=0}^\infty\frac{z^n}{n!}
 \int_\R x^ng(x)\,\mathrm d\lambda_1(x)=0
 \qquad(z\in\C).
\]
Dla $z=-i\xi$ klasyczna transformata całkowa $g$ znika.
Dla $\varphi\in\mathcal S$ z Fubiniego, z majorantą
$|g(x)\varphi(\xi)|$, mamy
\[
 \langle\mathcal F T_g,\varphi\rangle
 =(2\pi)^{-1/2}\int_\R\varphi(\xi)
 \int_\R g(x)e^{-ix\xi}\,\mathrm d\lambda_1(x)
 \,\mathrm d\lambda_1(\xi)=0.
\]
Twierdzenie~\ref{thm:fourier-auto-tempered22} i
Stwierdzenie~\ref{prop:schwartz-embeddings22} dają
$g=0$ prawie wszędzie, więc $f=0$. Przypadek rzeczywisty
wynika przez rozpatrzenie funkcji jako zespolonej.
W wymiarze $d$ ortonormalność iloczynów wynika z Fubiniego.
Zupełność wykazujemy indukcyjnie: jeśli $f$ jest
ortogonalna do wszystkich iloczynów, to dla każdego $n$
funkcja $g_n(x')=\int_\R f(x',t)h_n(t)\,\mathrm d\lambda_1(t)$
należy do $L^2(\R^{d-1})$ (Cauchy--Schwarz i Fubini)
i jest ortogonalna do wszystkich iloczynów w $d-1$
wymiarach. Zatem $g_n=0$ prawie wszędzie.
Przeliczalność rodziny $n$ pozwala wybrać wspólny zbiór
pełnej miary, na którym $f(x',\cdot)\in L^2(\R)$ i
wszystkie jej współczynniki znikają. Zupełność
jednowymiarowa i Fubini dają $f=0$.

Niech $L=-\Delta+|x|^2$ na $\mathcal S(\R^d)$.
Z rekurencji indukcyjnie dostajemy $P_n'=2nP_{n-1}$ dla $n\ge1$:
$P_{n+1}'=2P_n+2n(2xP_{n-1}-P_{n-1}')=2(n+1)P_n$,
z przypadkiem początkowym $P_1'=2P_0$.
Wprowadzamy $a_j=\partial_j+x_j$ i $a_j^*=-\partial_j+x_j$,
formalnie sprzężone na $\mathcal S$ przez całkowanie przez części.
Rekurencja, wzór na pochodną i normalizacja dają
\[
 a_jh_\nu=\sqrt{2\nu_j}\,h_{\nu-e_j},\qquad
 a_j^*h_\nu=\sqrt{2(\nu_j+1)}\,h_{\nu+e_j},
\]
gdzie pierwszy wyraz oznacza zero, gdy $\nu_j=0$.
Ponieważ $L=\sum_j a_j^*a_j+dI$, mamy
$Lh_\nu=(2|\nu|+d)h_\nu$.
Dla skończonej kombinacji $u=\sum_\nu c_\nu h_\nu$
i słowa $W$ długości $l$ z operatorów $a_j,a_j^*$
każdy $Wh_\nu$ jest zerem albo wielokrotnością
$h_{\nu+\gamma}$, gdzie $\gamma$ zależy wyłącznie od $W$.
Moduł współczynnika nie przekracza
$(2(|\nu|+l+1))^{l/2}$, ponieważ każdy krok zwiększa
stopień indeksu co najwyżej o jeden.
Niezerowe obrazy różnych indeksów są ortogonalne.
Dlatego
\[
 \|Wu\|_2^2\le C_l\sum_\nu(1+|\nu|)^l|c_\nu|^2.
\]
Z $x_j=(a_j+a_j^*)/2$ i $\partial_j=(a_j-a_j^*)/2$
oraz nierówności trójkąta wynika oszacowanie normy
każdego iloczynu $l$ operatorów mnożenia i różniczkowania
przez $C_l'\|(1+L)^r u\|_2$, gdy $2r\ge l$.

Dla całkowitego $s>d/2$ i $v\in\mathcal S$, ze wzoru
odwrotnego Fouriera (Twierdzenie~\ref{thm:fourier-schwartz17}),
nierówności Cauchy'ego--Schwarza i twierdzenia
Plancherela (Twierdzenie~\ref{thm:plancherel22}) wynika
\[
 \|v\|_\infty\le (2\pi)^{-d/2}\|\widehat v\|_1
 \le C_{d,s}\Big(\sum_{|\gamma|\le s}
 \|\partial^\gamma v\|_2^2\Big)^{1/2}.
\]
Stosujemy wagę $(1+|\xi|^2)^{-s/2}$, której kwadrat
jest całkowalny; rozwinięcie $(1+|\xi|^2)^s$ daje
skończoną sumę kwadratów monomów z dodatnimi współczynnikami.
Stosując to do $v=x^\alpha\partial^\beta u$ oraz
$(1+|x|)^N\le C_{N,d}\sum_{|\alpha|\le N}|x^\alpha|$,
dla $2r\ge N+m+s$ otrzymujemy
\[
 Q_{N,m}(u)\le C_{N,m,d,s}\|(1+L)^r u\|_2.
\]
Nierówność ta dotyczy na razie skończonych kombinacji.

Dla dowolnego $f\in\mathcal S$, rozwijając $(1+L)^r$
w skończoną sumę operatorów $x^\alpha\partial^\beta$
za pomocą reguły Leibniza, szacujemy każdy składnik
w $L^2$ przez stałą razy $Q_{|\alpha|+d+1,|\beta|}(f)$,
gdyż $(1+|x|)^{-d-1}\in L^2$.
Zatem $\|(1+L)^r f\|_2\le C_r Q_{N_r,m_r}(f)$.
Z formalnej symetrii $L$ na $\mathcal S$ mamy
$\langle(1+L)^r f,h_\nu\rangle
=(1+2|\nu|+d)^r\langle f,h_\nu\rangle$.
Parseval daje
\[
 \|(1+L)^r f\|_2^2
 =\sum_\nu(1+2|\nu|+d)^{2r}|H(f)_\nu|^2.
\]
Stąd $p_k(H(f))\le\|(1+L)^k f\|_2$, co dowodzi
ciągłości $H$ i przynależności współczynników do $s$.
Różnowartościowość wynika z zupełności bazy i ciągłości $f$.

Dla $c\in s(\N_0^d)$ i całkowitego $t$ z $2t>d$ mamy
\[
 \sum_\nu(1+2|\nu|+d)^{2r}|c_\nu|^2
 \le C_{r,d}p_{r+t}(c)^2
 \sum_\nu(1+|\nu|)^{-2t}<\infty
\]
na mocy zliczania wielowskaźników w dowodzie
Twierdzenia~\ref{thm:nuclear-s23}.
Sumy $u_M=\sum_{|\nu|\le M}c_\nu h_\nu$ są Cauchy'ego
w każdej normie $\|(1+L)^r\cdot\|_2$, a więc również
we wszystkich seminormach Schwartza. Zupełność
$\mathcal S$ daje granicę $u$. Ciągłość włożenia
$\mathcal S\to L^2$ (Stwierdzenie~\ref{prop:schwartz-operators17})
zapewnia $H(u)_\nu=c_\nu$ dla każdego $\nu$.
Otrzymujemy surjektywność $H$ i, przechodząc do granicy
w oszacowaniach sum częściowych,
$Q_{N,m}(H^{-1}c)\le C p_{r+t}(c)$.
Dowodzi to ciągłości odwrotności.
\end{proof}

\begin{wniosek}[Nuklearność przestrzeni Schwartza]
\label{cor:nuclear-schwartz23}
Przestrzeń Schwartza $\mathcal S(\R^d)$ jest nuklearną
przestrzenią Frécheta--Montela.
\end{wniosek}
\begin{proof}
Z Lematu~\ref{lem:hermite-schwartz23} mamy izomorfizm
liniowo-topologiczny $H:\mathcal S\to s(\N_0^d)$.
Dla ciągłej seminormy $p$ na $\mathcal S$ seminorma
$p\circ H^{-1}$ jest ciągła na $s(\N_0^d)$.
Z Twierdzenia~\ref{thm:nuclear-s23} wybieramy dominującą
seminormę $q'$ z nuklearnym przejściem między
uzupełnieniami. Po wzięciu $q=q'\circ H$ izometrie
indukowane przez $H$ na uzupełnieniach przenoszą
to przejście na $j_{q,p}$; własność ideałowa
(Stwierdzenie~\ref{prop:nuclear-ideal23}) daje jego
nuklearność. Zupełność, metryzowalność, beczkowość
i względna zwartość zbiorów ograniczonych są zachowywane
przez izomorfizmy liniowo-topologiczne, zatem również
własność Frécheta--Montela przechodzi na $\mathcal S$.
\end{proof}

\begin{uwaga}[Zakres twierdzenia o jądrze]
\label{rem:kernel-scope23}
Twierdzenie~\ref{thm:schwartz-kernel23} opisuje wszystkie
ciągłe operatory $\mathcal D(\Omega)\to\mathcal D'(\Lambda)_\beta$.
Nie oznacza ono, że jądro jest funkcją ani że operacja
składania dowolnych takich operatorów jest automatycznie
zdefiniowana. Twierdzenia trwałości nuklearności zostały udowodnione
powyżej; następny rozdział omawia topologiczne produkty tensorowe. Pełne rozwiązywanie równań
różniczkowych pozostaje w osobnym skrypcie.
\end{uwaga}
